\chapter{DishBrain}
\chaptermark{DishBrain}
\label{sec:dishbrain}

\section{접시 속 네트워크가 테니스를 친다}

DishBrain은 전극이 있는 칩 위에서 키운 살아 있는 뉴런 배양이 라켓 하나짜리 단순화된 테니스를 치는 테스트 벤치에 저자들이 붙인 이름이다~\cite{Kagan2022}. 살아 있는 뉴런으로 만든 기계를 다룬 실험 가운데 가장 많이 논의된 것이며(이런 기계의 개관은 \ref{sec:livingmachine}장), 이 책에는 특히 알맞다. 여기서는 작은 네트워크 전체에 손이 닿는다. 입력은 모두 실험자가 정하고, 출력은 모두 기록된다. 앞 장들의 질문, 즉 신호가 어떻게 부호화되는지, 누가 가르치는지, 프로브에 무엇이 보이는지를 눈도 근육도 \nt{DOPA} 뉴런도 없는 네트워크에 던질 수 있다. 엔지니어가 남의 테스트 벤치를 분석하듯 이 실험을 뜯어보자. 회로, 입력, 출력, 교사, 결과, 논쟁 순이다.

\section{테스트 벤치}

뉴런은 생쥐 태아의 피질에서 칩당 약 $800\,000$개를 얻거나, 인간 줄기세포에서 칩당 약 100만 개를 키운다. 뉴런은 몇 주 동안 칩 위에서 자라며 서로 얽혀 들어가고, 스스로 스파이크와 버스트를 내기 시작한다. 그 아래 $8$~mm$^2$ 면적에 백금 전극 $26\,000$개가 깔려 있다. 동시에 기록할 수 있는 것은 최대 $1024$개이고, 자극은 실제로 여덟 개를 통해 준다.

배양 둘레에 루프가 닫혀 있다(그림~\ref{fig:dishloop}). 컴퓨터가 게임을 모의한다. 공이 날아가 벽에 튕기고, 라켓이 위아래로 움직인다. 공의 위치는 여덟 개의 “감각” 전극에 흐르는 전류로 바뀐다. 두 “운동” 영역의 스파이크를 $10\ms$ 창마다 세어 라켓을 움직인다. 매번 공을 치거나 놓친 뒤 배양은 자극을 하나 더 받는다. 이것이 되먹임이다. $10\ms$ 창은 배양이 아니라 테스트 벤치 컴퓨터의 클록이다. 네트워크 자체는 이 책의 모든 네트워크처럼 여전히 비동기로 작동한다.

\begin{figure}[t]
\centering
\begin{tikzpicture}[>=Stealth,
  blk/.style={draw,very thick,rounded corners=3pt,align=center,font=\small,minimum height=1.2cm},
  lab/.style={font=\scriptsize,align=center}]
\node[blk,fill=blue!6,text width=3.4cm] (C) at (0,0) {칩 위의 배양\\\scriptsize 뉴런 $\sim10^6$개, 전극 $26\,000$개};
\node[blk,fill=orange!10,text width=3.0cm] (G) at (7.2,0) {컴퓨터:\\테니스 게임};
\draw[->,very thick,blue!60!black] (G.north west) to[bend right=25] node[lab,above] {공: \emph{장소} --- 8개 전극 중 어느 것,\\\emph{빈도} --- 초당 4--40회} (C.north east);
\draw[->,very thick,red!70!black] (C.south east) to[bend right=25] node[lab,below] {라켓: “위”와 “아래” 영역의\\$10\ms$당 스파이크 수} (G.south west);
\draw[->,thick,densely dashed,green!45!black] (G.east) -- ++(0.9,0) |- ++(0,-2.6) -| node[lab,pos=0.25,below] {히트: 예측 가능한 자극, 미스: 예측 불가능한 자극} (C.south);
\end{tikzpicture}
\caption{DishBrain 루프. 파란 화살표는 입력이다. 공의 위치를 위치 부호와 발화율 부호로 적는다(\ref{sec:code}장). 빨간 화살표는 출력이다. 두 영역의 스파이크 수 차이가 라켓을 움직인다. 초록 점선은 랠리마다 주는 되먹임으로, 테스트 벤치의 유일한 “교사”이다.}
\label{fig:dishloop}
\end{figure}

\section{입력과 출력}

\emph{입력}은 \ref{sec:code}장의 부호 두 가지로 동시에 적는다. 여덟 전극 가운데 어느 것을 자극하느냐가 공의 높이를 알려 준다. 이것은 \ref{sec:sensors}장의 센서처럼 위치 부호이다. 자극 펄스가 얼마나 자주 오느냐가 공이 얼마나 먼지를 알려 준다. 공이 먼 벽에 있으면 초당 $4$회이고, 라켓 쪽 벽에 가까워질수록 선형으로 늘어 초당 $40$회가 된다. 이것은 발화율 부호이며, 공 신호 펄스의 진폭은 $75$~mV이다. 배양은 이 부호들 가운데 어느 것도 미리 “알지” 못한다. 부호는 실험자가 정했고, 네트워크는 그것을 읽는 법을 배워야 한다.

\emph{출력}은 \ref{sec:debugging}장의 인터페이스와 같은 방식으로 읽는다. 한 운동 영역의 스파이크는 라켓을 위로, 다른 영역의 스파이크는 아래로 민다. 가장 단순하게 쓰면 창마다
\begin{empheq}[box=\keybox]{equation}
\Delta p \;=\; c\,\bigl(n_\uparrow - n_\downarrow\bigr),
\label{eq:dishreadout}
\end{empheq}
여기서 $n_\uparrow$, $n_\downarrow$는 $10\ms$ 동안 두 영역의 스파이크 수이고, $c$는 테스트 벤치의 계수이다. 이것은 \ref{sec:glutcomputer}장의 이중 레일 부호이다. 움직임의 부호는 신호의 부호가 아니라 두 전선 가운데 어느 쪽이 더 큰 소리를 내느냐가 전한다.

이런 출력의 잡음을 보자. 한 영역이 모두 합쳐 초당 $R$개의 스파이크를 낸다면, 창 $T=10\ms$ 동안 평균 $RT$개이고, \eqref{eq:rateerror}에 따라 상대 산포는 $1/\sqrt{RT}$이다. $R=1000$이면 창마다 약 $30$퍼센트, $R=100$이면 100퍼센트가 넘는다. 라켓은 떨릴 수밖에 없고, 네트워크가 배울 수 있는 방법은 하나뿐이다. 스파이크 수 차이의 \emph{평균}을 원하는 쪽으로 옮기는 것이다. \ref{sec:debugging}장의 모든 인터페이스를 제한하는 것과 같은 법칙이 이제는 루프의 양쪽에 걸린다.

\section{도파민 없는 교사}

이 테스트 벤치에서 가장 흥미로운 것은 교사이다. 피질 뉴런 배양에는 \nt{DOPA} 뉴런이 없고, 테스트 벤치는 “예상보다 좋다”는 신호를 전혀 보내지 않는다. 되먹임은 다른 것이다. 네트워크가 공을 받아 치면 여덟 감각 전극 모두에 짧고 \emph{예측 가능한} 자극이 한꺼번에 들어간다. $75$~mV 펄스를 초당 $100$회, $0.1\,\text{s}$ 동안이다. 놓치면 저자들이 \emph{예측 불가능하다}고 부르는 자극이 4초 동안 이어진다. $150$~mV 펄스를 초당 $5$회, 무작위 시점에 무작위로 고른 전극에 준다. 그 뒤 자극 없는 4초의 휴지가 있고, 공이 다시 출발한다~\cite{Kagan2022}.

저자들은 네트워크가 자기 입력이 예측 가능해지도록 행동한다는 가설에서 출발했다~\cite{Friston2010}. 엔지니어에게 의미는 간단하다. 배양이 4초의 무작위 잡음에서 벗어나는 방법은 정확히 하나, 공을 받아 치는 것이다. 같은 생각을 우리는 \ref{sec:code}장의 스파이크 절약과 \ref{sec:recognition}장의 프레임 예측에서 만났다.

그런데 여기서 꼭 예측 가능성이 필요할까? 더 거친 메커니즘도 가능하다. 미스 뒤의 예측 불가능한 자극이 네트워크의 최근 변화를 그저 \emph{뒤섞고}, 히트 뒤의 잔잔한 자극은 그것을 건드리지 않는다고 하자. 네트워크의 설정을 벡터 $\boldsymbol\psi$ 하나로 나타내고, 설정 $\boldsymbol\psi$에서 네트워크가 확률 $P_{\text{miss}}(\boldsymbol\psi)$로 공을 놓친다고 하자. 흔들기가 대칭이면, 즉 $\boldsymbol\psi$에서 $\boldsymbol\psi'$로 갈 확률이 되돌아올 확률과 같으면, 정상 상태에서 각 설정을 떠나는 흐름은 들어오는 흐름과 균형을 이루고, 네트워크가 설정 $\boldsymbol\psi$에 머무는 시간의 비율은
\begin{empheq}[box=\keybox]{equation}
\rho(\boldsymbol\psi) \;\propto\; \frac{1}{P_{\text{miss}}(\boldsymbol\psi)} .
\label{eq:dishdensity}
\end{empheq}
열 배 덜 놓치는 설정은 열 배 오래 유지된다. 가중치를 어느 쪽으로 바꿀지 네트워크에 알려 주는 것은 아무것도 없다. 좋은 설정이 덜 자주 부서질 뿐이다.

우리는 배양을 두 숫자로 줄인 모델로 이것을 확인했다. 공이 높이 $y$에 오면 라켓은 점 $a\,y+b$에 오고, 빗나간 거리가 $0.1$보다 작으면 받아 친다. 미스 뒤에는 두 숫자 모두 무작위로 흔들리고, 히트 뒤에는 그대로 남는다. 랠리 $2000$번 동안 받아 친 비율은 처음 200번의 $0.21$에서 마지막 500번의 $0.38$로 올랐다(“배양” 40개, 그림~\ref{fig:dishmodel}). 흔들기를 \emph{모든} 랠리 뒤에 주면 되먹임에 정보가 없고, 비율은 약 $0.12$에 머문다. 흔들기가 아예 없으면 $0.19$에 멈춰 있다. 테스트 벤치의 실험도 이와 맞는다. 세션 안의 유의미한 향상은 완전한 되먹임일 때만 있었다. 자극 대신 침묵을 주었을 때도 배양은 세션 끝에 되먹임이 없을 때보다 잘 쳤지만 효과가 작았다~\cite{Kagan2022}. 이 조건에서도 미스 뒤에는 긴 휴지가, 히트 뒤에는 짧은 휴지가 오므로 결과 사이의 차이가 완전히 사라지지는 않는다는 점에 주의하자.

\begin{figure}[t]
\centering
\begin{tikzpicture}[x=1cm,y=5cm]
\draw[->,gray!70!black] (0,0) -- (10.6,0);
\draw[->,gray!70!black] (0,0) -- (0,0.5) node[above,font=\small,black] {받아 친 비율};
\foreach \v in {0.1,0.2,0.3,0.4}{\draw[gray!60!black] (-0.06,\v) -- (0.06,\v); \node[left,font=\scriptsize] at (0,\v) {\v};}
\foreach \x/\f/\l/\lab in {1.8/0.21/0.38/{미스 뒤\\잡음},5.2/0.13/0.11/{모든 랠리 뒤\\잡음},8.6/0.19/0.19/{잡음 없음}}{
  \fill[blue!35] (\x-0.8,0) rectangle (\x-0.05,\f);
  \fill[red!60!black] (\x+0.05,0) rectangle (\x+0.8,\l);
  \node[below,font=\scriptsize,align=center] at (\x,-0.01) {\lab};
  \node[above,font=\scriptsize] at (\x-0.425,\f) {\f};
  \node[above,font=\scriptsize] at (\x+0.425,\l) {\l};}
\fill[blue!35] (7.4,0.47) rectangle (7.7,0.495); \node[right,font=\scriptsize] at (7.75,0.482) {처음 200};
\fill[red!60!black] (7.4,0.41) rectangle (7.7,0.435); \node[right,font=\scriptsize] at (7.75,0.422) {마지막 500};
\end{tikzpicture}
\caption{“미스 때 뒤섞기” 모델(\texttt{dishbrain\_model.py}): 랠리 $2000$번의 처음과 끝에서 받아 친 비율, 모델 배양 40개의 평균. 흔들기가 히트가 아니라 미스 뒤에 올 때만 학습이 일어난다. 세션 안의 유의미한 향상이 완전한 되먹임일 때만 있었던 실험과 같다.}
\label{fig:dishmodel}
\end{figure}

\begin{keyblock}
\begin{proposition}[뒤섞는 교사]
\label{prop:dishscramble}
실패가 네트워크를 흔들고 성공이 네트워크를 가만두면, 네트워크는 스스로 덜 틀리는 설정 쪽으로 표류한다. 한 설정에 머무는 시간은 실패 확률에 반비례한다~\eqref{eq:dishdensity}. 이런 학습에는 보상 신호도, 그 부호도, 예측 가능성 자체도 필요 없다. 성공 뒤의 되먹임과 실패 뒤의 되먹임이 다르기만 하면 된다. 배양의 행동 변화는 측정되었다. 접시 안에서 작동하는 메커니즘이 입력 예측인지 뒤섞기인지는 밝혀지지 않았다.
\end{proposition}
\end{keyblock}

\ref{sec:dofa}장과 비교해 보자. 거기서도 잡음은 탐색 역할을 했지만, 도파민에는 부호가 있었다. 오차 $\delta$가 가중치를 어느 쪽으로 옮길지 알려 주었고, 걸음은 경사를 따랐다~\eqref{eq:perturbation}. 여기에는 부호가 없고 “두기”와 “흔들기”만 있다. 이런 탐색은 더 거칠고 느리지만 네트워크에 요구하는 것이 적다. \nt{DOPA} 뉴런도, 비평가도, 몇 초에 걸친 흔적도 필요 없다.

\section{무엇을 측정했고 무엇을 두고 다투는가}

게임 한 판은 $20$분이었고, 처음 $5$분과 마지막 $15$분을 비교했다. 완전한 되먹임을 받은 배양에서는 랠리가 길어지고 즉시 미스가 줄었다. 세포 없는 대조, 게임 없는 대조, 컴퓨터 “라켓” 대조에서는 그런 일이 없었다. 인간 세포 배양은 평균적으로 생쥐 배양보다 오래 받아 쳤다. 향상은 통계적으로 확실하지만(생쥐 배양 9개와 인간 배양 11개, 각 그룹 100판 남짓) 작다. 네트워크는 좋은 선수가 되지 않았고, 무작위보다 조금 나아졌을 뿐이다. 향상이 확실한 것은 세션 안에서이고, 며칠에 걸친 세션 간의 안정적인 학습은 저자들도 얻지 못했다~\cite{Kagan2022}. 같은 그룹의 후속 연구는 게임 중 배양의 활동이 \emph{임계} 상태, 즉 감쇠와 사태 사이의 경계에 다가간다는 것을 발견했다. \ref{sec:gaba}장의 바로 그 경계 위의 삶이다. 가까울수록 게임을 더 잘했다~\cite{Habibollahi2023}.

가장 큰 논쟁을 부른 것은 숫자가 아니라 표현이었다. 논문 제목은 접시 속 뉴런이 “지각력”(sentience)을 보인다고 주장했고, 한 무리의 연구자들이 반박했다. 닫힌 루프 안의 행동 변화는 아직 지능도 감각도 아니며, 결론은 더 겸손하게 써야 한다는 것이다~\cite{Balci2023}. 공학적으로는 두 질문이 열려 있고, 둘 다 메커니즘에 관한 것이다. 첫째, 네트워크가 학습하는가, 즉 가중치가 오래 바뀌는가, 아니면 되먹임은 현재 상태를 옮길 뿐인가? 둘째, 꼭 예측 가능성이 필요한가, 아니면 명제~\ref{prop:dishscramble}에서처럼 성공과 실패에 대한 응답이 다르기만 하면 충분한가? 모든 전극에 손이 닿는 테스트 벤치는 두 질문 모두에 답할 수 있게 해 준다. 아마 이것이 이 테스트 벤치의 가장 큰 가치일 것이다.

\section{테스트 벤치 사양: 재현하는 법}
\label{sec:dishspec}

아래에는 이런 테스트 벤치를 다시 조립하는 데 필요한 모든 것을 모았다. 장비, 루프, 입력 부호화, 출력 판독, 되먹임, 실험 프로토콜이다. 매개변수마다 출처를 표시했다. “논문”은 값을 논문~\cite{Kagan2022}의 본문에서 가져왔다는 뜻이다. “선택”은 논문에 값이 없어서 재현하려면 직접 정해야 한다는 뜻이다. 이때는 기본값을 제안하고 그것을 확인하는 방법을 적는다.

\subsection*{장비와 배양}

\begin{center}
\footnotesize
\setlength{\tabcolsep}{4pt}
\renewcommand{\arraystretch}{1.15}
\begin{tabular}{>{\raggedright}p{4.0cm}>{\raggedright}p{6.9cm}>{\raggedright\arraybackslash}p{1.6cm}}
\toprule
매개변수 & 값 & 출처 \\
\midrule
칩 & MaxOne 배열: $8$~mm$^2$ 면적에 백금 전극 $26\,000$개 & 논문 \\
기록 & 동시에 최대 $1024$채널, 초당 $20\,000$ 표본 & 논문 \\
자극 & 기술적으로는 최대 $32$개 전극, 실험에서는 독립 전극 $8$개 & 논문 \\
세포 & 생쥐 태아 피질. 줄기세포에서 얻은 인간 뉴런(두 가지 배양법). 대조용 HEK293T 세포(뉴런 아님) & 논문 \\
파종 & 칩당 약 $10^6$개 세포 & 논문 \\
실험 시점의 배양 나이 & 생쥐: $\sim10$일부터. 인간: 배양법에 따라 $14$--$24$일 또는 $\sim80$일 & 논문 \\
\bottomrule
\end{tabular}
\end{center}

\subsection*{루프와 시간}

루프는 $10\ms$ 창($200$ 표본) 단위로 돈다. 창마다 운동 영역 전극의 스파이크를 세고, 게임을 갱신하고, 다음 자극을 내보낸다. 배양 속 스파이크에서 응답 자극까지의 지연은 약 $5\ms$이다. 게임 중 각 전극의 신호는 차단 주파수 $100$~Hz인 2차 베셀 고역 통과 필터를 지난다. 신호의 절댓값은 $1$~Hz의 1차 베셀 저역 통과 필터로 평활하며, 스파이크 문턱값은 이 평활된 절댓값에 비례한다. 배경의 표준편차 여섯 배라는 문턱값은 논문이 게임이 아니라 개별 활동 측정에 대해 밝힌 것이다. 자극이 배양의 응답으로 세어지지 않도록, $75$~mV를 넘는 큰 스파이크가 동시에 $15$개보다 많이 오면 모든 신호를 차단한다. 모두 논문에 따른 것이다.

\subsection*{공의 부호화}

\begin{center}
\footnotesize
\setlength{\tabcolsep}{4pt}
\renewcommand{\arraystretch}{1.15}
\begin{tabular}{>{\raggedright}p{3.4cm}>{\raggedright}p{7.5cm}>{\raggedright\arraybackslash}p{1.6cm}}
\toprule
매개변수 & 값 & 출처 \\
\midrule
감각 영역 & 전극 $8$개. 이웃한 전극이 이웃한 공 위치를 뜻하도록 배치 & 논문 \\
위치 부호 & 전극 번호 $k=1,\dots,8$이 라켓 중심에 대한 공의 위치를 정한다 & 논문 \\
어느 좌표인가 & 공의 수직 위치. 재현에서는 라켓 기준으로, $y-p\in[-\tfrac12,\tfrac12]$일 때 $k=\lceil 8(y-p+\tfrac12)\rceil$ & 논문; 공식은 선택 \\
발화율 부호 & 초당 $4$--$40$회의 자극 빈도가 라켓까지의 거리를 전한다 & 논문 \\
눈금 방향 & 공이 반대쪽 벽에 있으면 초당 $4$회, 라켓 쪽 벽에 가까워질수록 선형으로 초당 $40$회까지: $f=4+36\,(1-x/L)$, $x$는 라켓 쪽 벽까지의 거리, $L$은 코트 너비 & 논문 \\
펄스 모양 & 짧은 직사각형 2상 펄스: 먼저 양의 위상, 다음 음의 위상 & 논문 \\
진폭 & $75$~mV. 억제된 세포를 억지로 발화시키지 않도록 고른 값 & 논문 \\
위상 길이 & 밝히지 않음. 자극 시스템의 표준 설정을 쓰고 실험 중에는 바꾸지 않는다 & 선택 \\
\bottomrule
\end{tabular}
\end{center}

\subsection*{라켓 판독}

논문에는 운동 영역이 두 개 있다. 첫째 영역의 스파이크는 라켓을 위로, 둘째 영역의 스파이크는 아래로 움직인다. 세포 밀도와 활동의 차이를 반영하도록 두 영역의 기여에 가중치를 준다~\cite{Kagan2022}. 정확한 공식은 나와 있지 않다. 재현에서는 $10\ms$ 창마다 규칙~\eqref{eq:dishreadout} $\Delta p=c\,(n_\uparrow-n_\downarrow)$를 쓰고, 휴지 상태의 평균 활동으로 두 영역을 맞추는 가중치를 준다(선택). 즉 $n_\uparrow$와 $n_\downarrow$를 게임 전에 잰 각 영역의 창당 평균 스파이크 수로 나눈다. 계수 $c$는 평균 스파이크 수 하나만큼의 차이가 창마다 라켓을 코트 높이의 $1\%$만큼 옮기도록 고른다(선택). 그러면 한 영역이 계속 우세할 때 라켓은 $1$~s 동안 코트 전체를 가로지른다.

칩 위의 영역 배치는 논문 본문에 좌표로 주어지지 않았고 모식도만 있다. 재현에는 겹치지 않는 전극 그룹 세 개면 충분하다. 가운데에 감각 전극 여덟 개, 양옆에 기록 전극 그룹 하나씩, 하나는 “위”, 다른 하나는 “아래”이다(선택).

\subsection*{게임}

코트 크기, 라켓 길이, 공의 속도는 논문에 없다. 공이 일정한 속도로 날아가 벽에 튕긴다는 말뿐이다. 재현용(선택): 코트는 $1\times1$, 라켓은 오른쪽 벽에 길이 $0.2$(모델~\texttt{models/dishbrain\_model.py}에서처럼 $|y-p|<0.1$이면 히트), 공은 $1$~s 만에 코트를 가로지른다. 한 번도 받아 치지 못하고 놓친 랠리는 “에이스”, 연속 세 번보다 많이 받아 친 랠리는 “긴 랠리”로 센다(논문).

\subsection*{되먹임}

\begin{center}
\footnotesize
\setlength{\tabcolsep}{4pt}
\renewcommand{\arraystretch}{1.15}
\begin{tabular}{>{\raggedright}p{2.6cm}>{\raggedright}p{8.3cm}>{\raggedright\arraybackslash}p{1.6cm}}
\toprule
사건 & 주는 것 & 출처 \\
\midrule
히트 & 예측 가능한 자극: 감각 전극 $8$개 모두에 동시에, $75$~mV, 초당 $100$회, $100\ms$. 그동안 평소의 공 자극은 중단된다 & 논문 \\
미스 & 예측 불가능한 자극: $150$~mV, 초당 $5$회, $4$~s, 무작위 시점에 $8$개 가운데 무작위로 고른 전극에. 그 뒤 $4$~s 동안 자극 없음, 공은 무작위 방향으로 출발 & 논문 \\
무작위성의 법칙 & 밝히지 않음. 재현에서는 펄스 시점을 평균 빈도 초당 $5$회의 푸아송 과정으로 & 선택 \\
\bottomrule
\end{tabular}
\end{center}

\subsection*{프로토콜과 대조 실험}

세션 하나는 $20$~분이며, 처음 $5$~분과 마지막 $15$~분을 비교한다(논문). 결과 척도는 세 가지이다. 평균 랠리 길이, 에이스 비율, 긴 랠리 비율. 실험 조건(논문):
\begin{itemize}
\item \emph{자극}: 위에서 설명한 완전한 루프.
\item \emph{침묵}: 히트나 미스의 자극 대신 같은 시간 동안 자극을 전혀 주지 않고, 그 뒤 공이 무작위 방향으로 출발한다.
\item \emph{되먹임 없음}: 미스 뒤에 공은 그냥 튕겨 계속 날아가고, 공 위치 자극은 이어진다.
\item \emph{휴지}: 게임은 진행되고 라켓은 활동에 따라 움직이지만 자극은 전혀 없다.
\item \emph{대조}: 배지만 있는 칩. HEK293T 세포. 세포 대신 모의실험을 쓴 “컴퓨터 속 배양”.
\end{itemize}
본 실험의 표본 크기: 생쥐 배양 $9$개(세션 $101$번), 인간 배양 $11$개($138$), 배지만 있는 칩 $6$개($80$), 휴지 배양 $20$개($42$), 모의실험 $3$개($38$), 모두 $399$세션. 되먹임 실험에서는 $3$일 동안 하루 $3$세션씩 $486$세션, 조건은 “자극”, “침묵”, “되먹임 없음”($n=27$, $17$, $15$). 처음 $5$분에서 마지막 $15$분으로 평균 랠리 길이가 늘어난 것은 살아 있는 배양에서만이다. 되먹임 실험에서 시간에 따른 유의미한 향상은 “자극” 조건에서만 있었다. 세션 끝에 “침묵” 조건은 “되먹임 없음”보다 여전히 나았지만 효과가 작았고, 3일 동안 세션 간의 안정적인 학습은 없었다~\cite{Kagan2022}.

\subsection*{테스트 벤치의 주기, 단계별로}

테스트 벤치 컴퓨터는 $10\ms$마다 다음을 한다.
\begin{enumerate}
\item 지난 창 동안 두 운동 영역의 스파이크 수 $n_\uparrow$, $n_\downarrow$를 읽고, 각 영역의 휴지 상태 평균 스파이크 수로 정규화한다.
\item 라켓을 $\Delta p=c\,(n_\uparrow-n_\downarrow)$만큼 옮기고 코트 가장자리 안으로 제한한다.
\item 공을 한 걸음 옮기고, 위쪽, 아래쪽, 왼쪽 벽에서 튕긴다.
\item 공이 오른쪽 벽에 닿았으면: $|y-p|<0.1$일 때 히트, 랠리 수가 하나 늘고 히트 자극($100\ms$)을 준다. 아니면 미스, 랠리를 기록하고 미스 자극($4$~s)을 준 뒤 $4$~s 쉬고 공이 다시 출발한다.
\item 그렇지 않으면 공의 수직 위치로 전극 $k$를, 라켓 쪽 벽까지의 거리로 빈도 $f$를 골라 다음 창까지 전극 $k$에 빈도 $f$로 $75$~mV 2상 펄스를 준다.
\end{enumerate}
이것만으로 발표된 것과 호환되는 테스트 벤치를 조립하고 이 장의 핵심 질문을 확인할 수 있다. 배양을 가르치는 것이 예측 가능성인가, 아니면 히트와 미스에 대한 응답의 단순한 차이인가? 그러려면 논문의 세 조건에 논문에 없는 네 번째 조건을 더하는 것이 좋다. 미스 뒤에 예측 불가능한 자극과 같은 길이, 같은 에너지의 \emph{예측 가능한} 자극을 주는 것이다. 그래도 배양이 학습한다면 중요한 것은 예측 가능성이 아니라 차이이다.


\section{정리}

\begin{table}[t]
\centering
\footnotesize
\setlength{\tabcolsep}{4pt}
\renewcommand{\arraystretch}{1.15}
\begin{tabular}{>{\raggedright}p{2.1cm}>{\raggedright}p{4.2cm}>{\raggedright}p{1.9cm}>{\raggedright\arraybackslash}p{4.6cm}}
\toprule
테스트 벤치의 부분 & 구조 & 이 책의 장 & 알려진 것 \\
\midrule
입력 & 장소(8개 전극 중 어느 것)와 빈도(초당 4--40회) & \ref{sec:code}, \ref{sec:sensors} & 실험자가 정함 \\
출력 & $10\ms$ 동안 두 영역의 스파이크 수 차이~\eqref{eq:dishreadout} & \ref{sec:glutcomputer}, \ref{sec:debugging} & 실험자가 정함. 잡음은 \eqref{eq:rateerror}에 따름 \\
교사 & 히트 뒤 예측 가능한 자극, 미스 뒤 예측 불가능한 자극. 도파민 없음 & \ref{sec:dofa} & 세션 안의 유의미한 향상은 완전한 되먹임일 때만. 침묵일 때는 효과가 더 작음. 세션 간에는 불안정 \\
메커니즘 & 입력 예측 또는 미스 때 뒤섞기~\eqref{eq:dishdensity} & \ref{sec:dofa}, \ref{sec:recognition} & 밝혀지지 않음. 둘 다 결과를 설명함 \\
네트워크 상태 & 게임 중 임계 상태에 접근 & \ref{sec:gaba} & 게임 성적과의 관련은 측정됨 \\
“지각력” & 저자들의 주장 & --- & 보이지 않음. 반박됨 \\
\bottomrule
\end{tabular}
\caption{엔지니어의 눈으로 본 DishBrain.}
\label{tab:dishbrain}
\end{table}

DishBrain은 가장 단순한 형태의 살아 있는 뉴런 기계이다. 입력은 부호로 주어지고, 출력은 스파이크 수로 읽히며, 교사는 성공과 실패에 대한 응답의 차이이다(표~\ref{tab:dishbrain}). 눈도 근육도 도파민도 없는 네트워크가 아주 조금 게임을 배웠고, 그것만으로도 이 네트워크는 이 책의 생각들을 시험하는 테스트 벤치가 된다. 메커니즘이 예측이든, 뒤섞기든, 제3의 무엇이든, 답은 \ref{sec:debugging}장이 권하는 방식으로 찾게 될 것이다. 마침내 전체가 보이는 기계에서 프로브와 테스트 신호와 부분 끄기로 찾는 것이다.

\section{연습문제}

\begin{exercise}[\lvlA]
\label{ex:22:1}
\emph{스파이크 수를 얼마나 옮겨야 하는가.} 두 영역이 합쳐서 초당 $R_\uparrow+R_\downarrow=2000$개의 스파이크를 내고, 흐름은 푸아송이다. (a)~$R=1000$과 $R=100$일 때 한 영역의 $10\ms$당 스파이크 수의 상대 산포는 얼마인가? (b)~공이 $1$~s 동안 날아가는 사이에 변위~\eqref{eq:dishreadout}가 더해진다. 평균 변위가 산포의 두 배가 되는 가장 작은 $R_\uparrow-R_\downarrow$는 얼마인가?
\end{exercise}

\begin{exercise}[\lvlA]
\label{ex:22:2}
\emph{학습을 보려면 랠리가 몇 번 필요한가.} 랠리는 서로 독립이고 받아 친 비율은 이항 분포를 따른다. 받아 친 비율의 증가가 차이의 표준편차 두 배를 넘으려면 두 기간 각각에 랠리가 몇 번 필요한가? (a)~$0.20$에서 $0.25$로; (b)~모델에서처럼 $0.21$에서 $0.38$로.
\end{exercise}

\begin{exercise}[\lvlB]
\label{ex:22:3}
\emph{\eqref{eq:dishdensity}의 유도.} 설정 $\boldsymbol\psi$는 유한 집합 위를 움직인다. 미스(확률 $P(\boldsymbol\psi)>0$) 때 네트워크는 대칭 확률 $K(\boldsymbol\psi,\boldsymbol\psi')$로 $\boldsymbol\psi'$로 옮겨 가고, 히트 때는 그대로 남는다. (a)~$\rho\propto1/P$가 정상 분포임을 증명하라. (b)~$P=0.9$와 $P=0.1$인 두 설정이 있을 때 미스 비율은 얼마인가?
\end{exercise}

\begin{exercise}[\lvlB]
\label{ex:22:4}
\emph{모델의 흡수 영역.} 라켓은 $p=\min\bigl(1,\max(0,ay+b)\bigr)$에 오고, 공은 높이 $y\sim U(0,1)$에 오며, $|p-y|<h=0.1$이면 히트이다. (a)~$|b|<h$이고 $|a-1+b|<h$이면 네트워크가 모든 공을 받아 친다는 것을 증명하고, 이 영역의 넓이를 구하라. (b)~$a\sim U(-1,1)$, $b\sim U(0,1)$일 때 평균 히트 비율을 수치로 구하라.
\end{exercise}

\begin{exercise}[\lvlB]
\label{ex:22:5}
\emph{모의실험: 설정 주위의 울타리.} \texttt{dishbrain\_model.py}의 사본에서 배양 $200$개에 각각 랠리 $20\,000$번을 주라. 마지막 $500$번의 히트 비율은 얼마인가? 흔들기마다 설정을 직사각형 $a\in[-1,2]$, $b\in[-1,1]$ 안으로 잘라 넣으면 어떻게 되는가? 연습문제~\ref{ex:22:3}을 써서 차이를 설명하라.
\end{exercise}

\begin{exercise}[\lvlB]
\label{ex:22:6}
\emph{입력 부호 설계.} 높이 오차가 $h=0.1$보다 작으면 공을 받아 친다. 네트워크는 $T=0.25$~s 동안 입력을 읽고, 빈도 $r$까지 약 $\sqrt{rT}$개의 수준을 구별하며, 자극은 초당 $40$회를 넘지 않는다. (a)~전극 여덟 개의 위치 부호로 충분한가? (b)~전극 하나의 빈도로 높이를 전할 수 있는가?
\end{exercise}

\begin{exercise}[\lvlC]
\label{ex:22:7}
\emph{연구: 예측인가 뒤섞기인가?} 모델을 조정 가능한 숫자 $d$개($p=\sum_{i<d}a_iy^i$)로 일반화하라. 히트 비율 $0.5$에 이르기까지 필요한 랠리 수는 뒤섞기일 때와 오차 부호 규칙~\eqref{eq:perturbation}일 때 $d$에 따라 어떻게 늘어나는가? 어떤 테스트 벤치 실험이 두 가설을 가려내겠는가?
\end{exercise}
